% CV, español. Compila las dos versiones con `npm run cv`.
\documentclass[a4paper,10pt]{article}
\usepackage[spanish,es-noshorthands,es-nolists,es-nodecimaldot]{babel}
% Shared layout for cv_en.tex and cv_es.tex.
% Base template: Jake Gutierrez's resume (MIT).
% Each CV loads its own babel language before % Shared layout for cv_en.tex and cv_es.tex.
% Base template: Jake Gutierrez's resume (MIT).
% Each CV loads its own babel language before % Shared layout for cv_en.tex and cv_es.tex.
% Base template: Jake Gutierrez's resume (MIT).
% Each CV loads its own babel language before \input{cv-style}.

\usepackage{latexsym}
\usepackage[empty]{fullpage}
\usepackage{titlesec}
\usepackage{marvosym}
\usepackage[usenames,dvipsnames]{color}
\usepackage{xcolor}
\usepackage{verbatim}
\usepackage{enumitem}
\usepackage[hidelinks]{hyperref}
\usepackage{fancyhdr}
\usepackage{tabularx}
\usepackage[utf8]{inputenc}
\usepackage[T1]{fontenc}

\input{glyphtounicode}
\pdfgentounicode=1

% No hyphenation: ATS parsers read "valida-" + "tion" as two broken words and
% miss the keyword. Lines may stretch a little instead.
\hyphenpenalty=10000
\exhyphenpenalty=10000
\emergencystretch=1.5em

%--- Page style ---
\pagestyle{fancy}
\fancyhf{}
\fancyfoot{}
\renewcommand{\headrulewidth}{0pt}
\renewcommand{\footrulewidth}{0pt}

%--- Margins ---
\addtolength{\oddsidemargin}{-0.65in}
\addtolength{\evensidemargin}{-0.65in}
\addtolength{\textwidth}{1.3in}
\addtolength{\topmargin}{-0.35in}
\addtolength{\textheight}{1.2in}

%--- Global spacing reset ---
\urlstyle{same}
\raggedbottom
\setlength{\tabcolsep}{0in}
\setlength{\parskip}{0pt}
\setlength{\parsep}{0pt}

%--- Section title formatting ---
\titleformat{\section}{
  \vspace{-6pt}\scshape\raggedright\large
}{}{0em}{}[\color{black}\titlerule\vspace{-4pt}]

%-------------------------
% Custom commands

\newcommand{\resumeItem}[1]{%
  \item\small{\linespread{0.8}\selectfont #1}\vspace{-2pt}%
}

\newcommand{\resumeSubheading}[4]{%
  \vspace{3pt}\item[]%
    \begin{tabular*}{\textwidth}[t]{@{}l@{\extracolsep{\fill}}r@{}}
      \textbf{#1} & \small #2 \\
      \textit{\small #3} & \textit{\small #4} \\
    \end{tabular*}\vspace{1pt}%
}

\newcommand{\jobdivider}{%
  \par\vspace{2pt}%
  \noindent\textcolor{gray!50}{\rule{\textwidth}{0.25pt}}%
  \vspace{-2pt}%
}

\newcommand{\resumeSubHeadingListStart}{%
  \begin{itemize}[leftmargin=0pt, label={}, topsep=0pt, itemsep=0pt, parsep=0pt]%
}
\newcommand{\resumeSubHeadingListEnd}{\end{itemize}}

\newcommand{\resumeItemListStart}{%
  \begin{itemize}[leftmargin=0.2in, topsep=1pt, itemsep=2pt, parsep=0pt]%
}
\newcommand{\resumeItemListEnd}{\end{itemize}\vspace{1pt}}

% Project heading: name and stack on the left, link on the right.
\newcommand{\resumeProjectHeading}[2]{%
  \vspace{3pt}\item[]%
    \begin{tabular*}{\textwidth}[t]{@{}l@{\extracolsep{\fill}}r@{}}
      \small #1 & \small #2 \\
    \end{tabular*}\vspace{1pt}%
}

% Private build (npm run cv:private) defines \cvprivate and adds the phone
% number from cv/private.tex, which git ignores. The public build
% (npm run cv) never reads that file, so the PDFs on the website have no phone.
\ifdefined\cvprivate
  \IfFileExists{private.tex}{\input{private.tex}}{%
    \errmessage{Missing cv/private.tex: copy cv/private.example.tex and put your phone number in it}}
\fi
\newcommand{\cvsep}{\hspace{0.3em}\textbar\hspace{0.3em}}
\newcommand{\cvphoneitem}{\ifdefined\cvphone\cvphone\cvsep\fi}
.

\usepackage{latexsym}
\usepackage[empty]{fullpage}
\usepackage{titlesec}
\usepackage{marvosym}
\usepackage[usenames,dvipsnames]{color}
\usepackage{xcolor}
\usepackage{verbatim}
\usepackage{enumitem}
\usepackage[hidelinks]{hyperref}
\usepackage{fancyhdr}
\usepackage{tabularx}
\usepackage[utf8]{inputenc}
\usepackage[T1]{fontenc}

\input{glyphtounicode}
\pdfgentounicode=1

% No hyphenation: ATS parsers read "valida-" + "tion" as two broken words and
% miss the keyword. Lines may stretch a little instead.
\hyphenpenalty=10000
\exhyphenpenalty=10000
\emergencystretch=1.5em

%--- Page style ---
\pagestyle{fancy}
\fancyhf{}
\fancyfoot{}
\renewcommand{\headrulewidth}{0pt}
\renewcommand{\footrulewidth}{0pt}

%--- Margins ---
\addtolength{\oddsidemargin}{-0.65in}
\addtolength{\evensidemargin}{-0.65in}
\addtolength{\textwidth}{1.3in}
\addtolength{\topmargin}{-0.35in}
\addtolength{\textheight}{1.2in}

%--- Global spacing reset ---
\urlstyle{same}
\raggedbottom
\setlength{\tabcolsep}{0in}
\setlength{\parskip}{0pt}
\setlength{\parsep}{0pt}

%--- Section title formatting ---
\titleformat{\section}{
  \vspace{-6pt}\scshape\raggedright\large
}{}{0em}{}[\color{black}\titlerule\vspace{-4pt}]

%-------------------------
% Custom commands

\newcommand{\resumeItem}[1]{%
  \item\small{\linespread{0.8}\selectfont #1}\vspace{-2pt}%
}

\newcommand{\resumeSubheading}[4]{%
  \vspace{3pt}\item[]%
    \begin{tabular*}{\textwidth}[t]{@{}l@{\extracolsep{\fill}}r@{}}
      \textbf{#1} & \small #2 \\
      \textit{\small #3} & \textit{\small #4} \\
    \end{tabular*}\vspace{1pt}%
}

\newcommand{\jobdivider}{%
  \par\vspace{2pt}%
  \noindent\textcolor{gray!50}{\rule{\textwidth}{0.25pt}}%
  \vspace{-2pt}%
}

\newcommand{\resumeSubHeadingListStart}{%
  \begin{itemize}[leftmargin=0pt, label={}, topsep=0pt, itemsep=0pt, parsep=0pt]%
}
\newcommand{\resumeSubHeadingListEnd}{\end{itemize}}

\newcommand{\resumeItemListStart}{%
  \begin{itemize}[leftmargin=0.2in, topsep=1pt, itemsep=2pt, parsep=0pt]%
}
\newcommand{\resumeItemListEnd}{\end{itemize}\vspace{1pt}}

% Project heading: name and stack on the left, link on the right.
\newcommand{\resumeProjectHeading}[2]{%
  \vspace{3pt}\item[]%
    \begin{tabular*}{\textwidth}[t]{@{}l@{\extracolsep{\fill}}r@{}}
      \small #1 & \small #2 \\
    \end{tabular*}\vspace{1pt}%
}

% Private build (npm run cv:private) defines \cvprivate and adds the phone
% number from cv/private.tex, which git ignores. The public build
% (npm run cv) never reads that file, so the PDFs on the website have no phone.
\ifdefined\cvprivate
  \IfFileExists{private.tex}{\input{private.tex}}{%
    \errmessage{Missing cv/private.tex: copy cv/private.example.tex and put your phone number in it}}
\fi
\newcommand{\cvsep}{\hspace{0.3em}\textbar\hspace{0.3em}}
\newcommand{\cvphoneitem}{\ifdefined\cvphone\cvphone\cvsep\fi}
.

\usepackage{latexsym}
\usepackage[empty]{fullpage}
\usepackage{titlesec}
\usepackage{marvosym}
\usepackage[usenames,dvipsnames]{color}
\usepackage{xcolor}
\usepackage{verbatim}
\usepackage{enumitem}
\usepackage[hidelinks]{hyperref}
\usepackage{fancyhdr}
\usepackage{tabularx}
\usepackage[utf8]{inputenc}
\usepackage[T1]{fontenc}

\input{glyphtounicode}
\pdfgentounicode=1

% No hyphenation: ATS parsers read "valida-" + "tion" as two broken words and
% miss the keyword. Lines may stretch a little instead.
\hyphenpenalty=10000
\exhyphenpenalty=10000
\emergencystretch=1.5em

%--- Page style ---
\pagestyle{fancy}
\fancyhf{}
\fancyfoot{}
\renewcommand{\headrulewidth}{0pt}
\renewcommand{\footrulewidth}{0pt}

%--- Margins ---
\addtolength{\oddsidemargin}{-0.65in}
\addtolength{\evensidemargin}{-0.65in}
\addtolength{\textwidth}{1.3in}
\addtolength{\topmargin}{-0.35in}
\addtolength{\textheight}{1.2in}

%--- Global spacing reset ---
\urlstyle{same}
\raggedbottom
\setlength{\tabcolsep}{0in}
\setlength{\parskip}{0pt}
\setlength{\parsep}{0pt}

%--- Section title formatting ---
\titleformat{\section}{
  \vspace{-6pt}\scshape\raggedright\large
}{}{0em}{}[\color{black}\titlerule\vspace{-4pt}]

%-------------------------
% Custom commands

\newcommand{\resumeItem}[1]{%
  \item\small{\linespread{0.8}\selectfont #1}\vspace{-2pt}%
}

\newcommand{\resumeSubheading}[4]{%
  \vspace{3pt}\item[]%
    \begin{tabular*}{\textwidth}[t]{@{}l@{\extracolsep{\fill}}r@{}}
      \textbf{#1} & \small #2 \\
      \textit{\small #3} & \textit{\small #4} \\
    \end{tabular*}\vspace{1pt}%
}

\newcommand{\jobdivider}{%
  \par\vspace{2pt}%
  \noindent\textcolor{gray!50}{\rule{\textwidth}{0.25pt}}%
  \vspace{-2pt}%
}

\newcommand{\resumeSubHeadingListStart}{%
  \begin{itemize}[leftmargin=0pt, label={}, topsep=0pt, itemsep=0pt, parsep=0pt]%
}
\newcommand{\resumeSubHeadingListEnd}{\end{itemize}}

\newcommand{\resumeItemListStart}{%
  \begin{itemize}[leftmargin=0.2in, topsep=1pt, itemsep=2pt, parsep=0pt]%
}
\newcommand{\resumeItemListEnd}{\end{itemize}\vspace{1pt}}

% Project heading: name and stack on the left, link on the right.
\newcommand{\resumeProjectHeading}[2]{%
  \vspace{3pt}\item[]%
    \begin{tabular*}{\textwidth}[t]{@{}l@{\extracolsep{\fill}}r@{}}
      \small #1 & \small #2 \\
    \end{tabular*}\vspace{1pt}%
}

% Private build (npm run cv:private) defines \cvprivate and adds the phone
% number from cv/private.tex, which git ignores. The public build
% (npm run cv) never reads that file, so the PDFs on the website have no phone.
\ifdefined\cvprivate
  \IfFileExists{private.tex}{\input{private.tex}}{%
    \errmessage{Missing cv/private.tex: copy cv/private.example.tex and put your phone number in it}}
\fi
\newcommand{\cvsep}{\hspace{0.3em}\textbar\hspace{0.3em}}
\newcommand{\cvphoneitem}{\ifdefined\cvphone\cvphone\cvsep\fi}

\hypersetup{pdftitle={Jesús Jiménez Alcántara, Ingeniero Backend, CV}, pdfauthor={Jesús Jiménez Alcántara}}

\begin{document}

%----------CABECERA----------
\begin{center}
    \textbf{\Huge \scshape Jesús Jiménez Alcántara} \\ \vspace{2pt}
    {\large Ingeniero Backend y Arquitecto de Soluciones Cloud}\hspace{0.6em}{\textbar}\hspace{0.6em}{\normalsize .NET $\cdot$ Azure $\cdot$ DevOps $\cdot$ IA} \\ \vspace{3pt}
    {\footnotesize Málaga, España\cvsep\cvphoneitem\href{mailto:jjalcantara.dev@gmail.com}{\underline{jjalcantara.dev@gmail.com}}\cvsep\href{https://linkedin.com/in/jjalcantara}{\underline{linkedin.com/in/jjalcantara}}\cvsep\href{https://github.com/jjalcantara-dev}{\underline{github.com/jjalcantara-dev}}\cvsep\href{https://jjalcantara.dev}{\underline{jjalcantara.dev}}}
\end{center}

\vspace{-6pt}

%-----------RESUMEN-----------
\section{Resumen Profesional}
\small
\noindent Ingeniero backend y arquitecto de soluciones cloud con más de 4 años diseñando y llevando a producción sistemas en
\textbf{C\#/.NET} sobre \textbf{Azure}: APIs REST, procesamiento orientado a eventos, pagos (Stripe, App Store, Google Play) y
seguridad (JWT, OAuth 2.0, RBAC), del modelo de datos a CI/CD y soporte en producción. Experiencia en ingeniería de seguridad
en Ernst \& Young; cursando un Máster en IA, con prácticas en herramientas para asistentes de programación con IA.

\vspace{-2pt}

%-----------EXPERIENCIA-----------
\section{Experiencia}
  \resumeSubHeadingListStart

    \resumeSubheading
      {Kaizten Analytics}{Remoto}
      {Ingeniero de IA, prácticas curriculares (Máster en IA)}{Septiembre 2026 -- Noviembre 2026}
    \resumeItemListStart
        \resumeItem{Desarrollé en solitario, en \textbf{Java 17}, las herramientas que conectan el analizador de Kaizten con Claude Code y Codex vía \textbf{MCP}.}
        \resumeItem{Publiqué \textbf{17 versiones} multiplataforma (instaladores para Linux, macOS y Windows, plugin de Maven, control de CI).}
        \resumeItem{Aseguré la calidad con \textbf{más de 1.500 tests automatizados} y documentación bilingüe (ES/EN) generada en cada versión.}
    \resumeItemListEnd

    \jobdivider

    \resumeSubheading
      {The Bubble Hub}{Málaga, España}
      {Ingeniero Backend}{Abril 2024 -- Abril 2026}
    \resumeItemListStart
        \resumeItem{Diseñé la arquitectura backend de dos apps móviles en \textbf{ASP.NET Core} sobre \textbf{Azure} (Clean Architecture, SOLID, Docker), que crecieron de decenas a \textbf{más de 3.000 usuarios activos mensuales}, y de la plataforma formativa DailyDermaTalks (Next.js).}
        \resumeItem{Me encargué de las APIs REST en producción de principio a fin: acordé contratos de datos y semántica de errores con los equipos de iOS y Android y versioné las especificaciones \textbf{OpenAPI} para que las apps se publicaran sin depender del backend.}
        \resumeItem{Moví las tareas largas fuera del ciclo petición-respuesta con \textbf{Azure Storage Queues} y Worker Services y eliminé la pérdida de webhooks de suscripción con \textbf{procesamiento idempotente} y colas dead-letter, a \textbf{decenas de miles de peticiones al día}.}
        \resumeItem{Construí la capa de pagos completa: \textbf{Stripe} (webhooks, claves de idempotencia y conciliación) y suscripciones de App Store y Google Play, con renovaciones, periodos de gracia, reembolsos y revocaciones.}
        \resumeItem{Diseñé el modelo de datos en \textbf{PostgreSQL} (normalización, índices, integridad referencial, migraciones) y la seguridad por capas: \textbf{JWT} stateless, RBAC, OAuth 2.0 (Google, Apple) y rate limiting con Redis.}
        \resumeItem{Sustituí los despliegues manuales por un pipeline de \textbf{Azure DevOps} con controles de calidad (tests automatizados, SonarCloud) y promoción entre desarrollo, preproducción y producción, con rollback y monitorización con Application Insights.}
    \resumeItemListEnd

    \jobdivider

    \resumeSubheading
      {Ernst \& Young (EY)}{Málaga, España}
      {Consultor de Tecnología y Seguridad}{Abril 2022 -- Abril 2024}
    \resumeItemListStart
        \resumeItem{Desarrollé componentes backend en ASP.NET para clientes enterprise regulados en entornos SecDevOps, con seguridad desde el diseño: validación de entradas, cifrado de datos y control de acceso por mínimo privilegio.}
        \resumeItem{Construí \textbf{pipelines de ingesta de logs de seguridad} que integraban \textbf{más de 10 fuentes heterogéneas} en IBM QRadar, Splunk y Wazuh, y desarrollé la capa de detección (casos de uso, búsquedas de correlación y playbooks de alertas) para el SOC.}
        \resumeItem{Optimicé CI/CD con Jenkins y GitHub Actions, con tests automatizados y controles de calidad bajo una gestión del cambio estricta.}
    \resumeItemListEnd

  \resumeSubHeadingListEnd

\vspace{-4pt}

%-----------OPEN SOURCE-----------
\section{Open Source}
  \resumeSubHeadingListStart
    \resumeProjectHeading
      {\textbf{StoreBridge} $|$ \emph{C\#, .NET 8/9/10, NuGet}}
      {\href{https://github.com/jjalcantara-dev/StoreBridge}{\underline{github.com/jjalcantara-dev/StoreBridge}}}
    \resumeItemListStart
        \resumeItem{Autor de una librería en \textbf{NuGet} para verificar desde el backend suscripciones y compras in-app de App Store y Google Play.}
        \resumeItem{Autentica criptográficamente los webhooks de las tiendas (cadena x5c de Apple anclada a Apple Root CA G3, tokens OIDC de Google Pub/Sub), con reintentos, validación de la configuración al arrancar y \textbf{197 tests unitarios}.}
    \resumeItemListEnd
  \resumeSubHeadingListEnd

\vspace{-4pt}

%-----------FORMACIÓN-----------
\section{Formación}
  \resumeSubHeadingListStart
    \resumeSubheading
      {Universidad Internacional de La Rioja (UNIR)}{La Rioja, España}
      {Máster en Inteligencia Artificial}{Previsto marzo 2027}
    \resumeSubheading
      {Universidad de Málaga (UMA)}{Málaga, España}
      {Grado en Ingeniería del Software}{Octubre 2022}
  \resumeSubHeadingListEnd

\vspace{-4pt}

%-----------HABILIDADES-----------
\section{Habilidades Técnicas}
\vspace{2pt}
\small
\begin{itemize}[leftmargin=0pt, label={}, topsep=0pt, itemsep=1pt, parsep=0pt]
  \item \textbf{Lenguajes y frameworks:} C\#, Java, Python, TypeScript, SQL; ASP.NET Core, Entity Framework Core, Spring Boot, Next.js
  \item \textbf{Cloud y DevOps:} Azure (Functions, DevOps, App Service, Storage), Docker, GitHub Actions, Jenkins, Maven, SonarCloud
  \item \textbf{Datos y mensajería:} PostgreSQL, SQL Server, CosmosDB, Redis, Azure Storage Queues, Worker Services
  \item \textbf{APIs y seguridad:} REST, OpenAPI/Swagger, JWT, OAuth 2.0, RBAC, rate limiting, Stripe, pagos de App Store/Google Play
  \item \textbf{IA, testing y prácticas:} Model Context Protocol (MCP), TDD, xUnit, SOLID, Clean Architecture, Agile/Scrum
  \item \textbf{Idiomas:} español (nativo), inglés (B2, competencia profesional)
\end{itemize}

\vspace{-4pt}

%-----------OTRA EXPERIENCIA-----------
\section{Otra Experiencia}
\small
\noindent\textbf{Junta de Andalucía}, San Pedro de Alcántara, España: profesor interino de Informática, junio 2026.

\end{document}
